\section{Related work}
\label{sec:related}

\paragraph{Reconstruction and local coordinates.}
Kernel affine hull machines (KAHMs) were introduced for geometrically informed
classification, including a differential-privacy application
\citep{Kumar2024KAHM}. We adapt their reconstruction operators to form
unsupervised gates on forecasting histories and compare them with locality
and affine-projection controls. Local coordinate coding already connects
localized representations to approximation of smooth prediction functions
\citep{Yu2009LCC}. Local tangent coding augments coding weights with weighted
displacements along local tangent directions \citep{YuZhang2010LTC}; our gated
local coordinates use this established architectural idea. The empirical
question is whether KAHM-derived gates improve forecasting within the stated
architecture and resource regimes. Classification results for the original
KAHMs and approximation results for local coding concern their own tasks
and assumptions; they do not establish the outcome of this comparison.

\paragraph{Target relevance and finite prediction procedures.}
Supervised dimension reduction explicitly seeks representations that preserve
the relationship between an input and its response \citep{Fukumizu2004SDR}.
This highlights why a representation's reconstruction quality alone does
not determine its forecasting value. Our raw-input augmentation retains the
original history, while changing the coordinates available to a particular
fitted predictor. Its finite-sample performance can therefore change through
approximation, estimation or regularization without a change in the
information contained in that history. We use finite Nystr\"om kernel
features \citep{Williams2001Nystrom}, so the comparison concerns specified
fitting procedures. Equal requested landmark counts after different feature
maps do not imply identical prediction classes, realized ranks or computation.
The capacity caps and measured costs make these distinctions explicit.

\paragraph{Forecast selection and performance prediction.}
Learning which forecasting method suits a series is an established problem.
\citet{Prudencio2004MetaLearning} studied both learned selection between
forecasting models and model ranking. FFORMA learns forecast-combination
weights from time-series features \citep{MonteroManso2020FFORMA}.
FFORMPP predicts forecasting errors from historical-series features and
uses those predictions for selection or combination
\citep{Talagala2022FFORMPP}; FFORMS learns a feature-based mapping to a
selected forecasting model \citep{Talagala2023FFORMS}. These works also
motivate attention to computation and the representativeness of the data
used to train a selector. Here a deliberately small, fixed ridge diagnostic
predicts the relative benefit of appending KAHM local features to a raw-history
kernel predictor. The main question is the incremental value of geometric
features beyond metadata and a fixed validation probe. Performance comparisons
apply to the implemented candidates and selectors. Four development families
and four reserved families support a bounded examination of this question,
with substantial limits on transfer to additional domains.

\paragraph{Cheap probes and the cost of choosing a method.}
Landmarking characterizes a learning task using the measured performance of
inexpensive learners \citep{Pfahringer2000Landmarking}. Our cheap validation
probe follows this general principle: it supplies evidence about the target
that input-only geometric features do not directly measure. We retain both
metadata-only and probe-augmented diagnostic controls, and charge the probe
and geometry computation even when the eventual decision selects the raw
predictor. The accounting also charges all searches needed by a validation
selector, followed by the chosen candidate's refit and forecast. Offline
diagnostic fitting and per-task inference are reported separately. This
evaluates the computational consequences of the declared decision rules;
the learned rule, the use of probes and the general accuracy--cost tradeoff
are established ideas.

\paragraph{Forecast evaluation and source separation.}
Reusing a finite model-selection criterion for evaluation can produce
optimistic conclusions \citep{Cawley2010Selection}. Forecasting adds concerns
about temporal information, scaling, appropriate simple baselines and
dependent errors \citep{Hewamalage2023Evaluation}. We use chronological
training, validation and test blocks; the separate development/confirmation
division concerns the original data sources from which those tasks arise.
The Monash archive supplies datasets from varied domains and documented
baseline evaluations \citep{Godahewa2021Monash}. We use pinned snapshots and
keep derived series, horizons, capacities and seeds within their source
families for aggregation. Diagnostic features and regularization are fixed
before development outcomes, and fitted diagnostics are frozen before
confirmation evaluation. The small, deliberately chosen family collection
supports descriptive source-level comparisons, not population-wide
confidence claims. This is a historical forecasting study; it does not
establish real-time deployment performance or performance under future
changes in data availability.
